\BourbakiStarChapter{记号索引}

\(x + y, x.y, xy, x \top y, x \bot y: I, \S 1, no. 1.\) \(X \top Y, X + Y, XY (X, Y subsets): I, \S 1, no. 1.\) \(X \top a, a \top X (X a subset, a an element): I, \S 1, no. 1.\) \(\prod_{\alpha \in A} x_{\alpha}, \prod_{\alpha} x_{\alpha}, \top x_{\alpha}, \underset{\alpha \in A}{\bot} x, \underset{\alpha}{\bot} x_{\alpha}, \bot x_{\alpha}, \sum_{\alpha \in A} x_{\alpha}, \sum_{\alpha} x_{\alpha}, \sum x_{\alpha}, \prod_{\alpha \in A} x_{\alpha}, \prod_{\alpha} x_{\alpha}, \prod_{\alpha} x_{\alpha}, \prod_{\alpha} x_{\alpha}\) \(\prod_{p \leqslant i \leqslant q} x_i, \prod_{i=p}^q x_i: I, \S 1, no. 2.\) \(x_p \top x_{p+1} \top \cdots \top x_q: I, \S 1, no. 3.\) \(\frac{n}{\top} x, \underset{i < j < n}{\bot} x^n, nx (n \in N): I, \S 1, no. 3.\) \(\sum_{i=p}^{q} \sum_{j=r}^{s} x_{ij}, \sum_{j=r}^{s} \sum_{i=p}^{q} x_{ij}: I, \S 1, no. 5.\) \(\sum_{0 \leqslant i_1 < i_2 < \cdots < i_p \leqslant n} ^{q} \prod_{i_1 i_2 \cdots i_p, i_1 < i_2 < \cdots < i_p} ^{q} x_{i_1 i_2} x_{i_1 i_2} x_{i_1 i_2} x_{i_1 i_2} x_{i_1 i_2} x_{i_1 i_2} x_{i_1 i_2} x_{i_1 i_2} x_{i_1 i_2} x_{i_1 i_2} x_{i_1 i_2} x_{i_1 i_2}\) \(0, 1: I, \S 2, no. 1.\) \(\gamma_a, \delta_a, \gamma(a), \delta(a): I, \S 2, no. 2.\) \(E_S (S a subset of a commutative monoid E): I, \S 2, no. 4.\) \(Z, +(addition in Z): I, \S 2, no. 5.\) \(\leqslant (order relation on Z): I, \S 2, no. 5.\) \(N^*: I, \S 2, no. 5.\) \(\frac{n}{\top} (for n \in Z): I, \S 2, no. 7.\) \(-x, x - y, x + y - z, x - y - z, x - y + z - t: I, \S 2, no. 8.\) \(nx (n \in Z): I, \S 2, no. 8.\) \(x^n (n \in Z): I, \S 2, no. 8.\) \(\frac{1}{x}, \frac{x}{y}, x/y: I, \S 2, no. 8.\)

\(\alpha \perp x, \alpha \perp X, \Xi \perp X\) ( \(\alpha\) 一个算子， \(\Xi\) 一组算子）：I，§3，第1节。 \(\mathfrak{S}_{\mathrm{F}}\) ：I，§4，第1节。

\(\alpha .x,x.\alpha ,x^{\alpha}\) ( \(\alpha\) 一个算子）：I，§3，第1节。

(G:H)，G/H（H为G的子群）：I，§4，第4节。

\(x \equiv y \pmod{H}, x \equiv y (H) (H a normal subgroup): I, § 4, no. 4.\)

\(\operatorname{Ker} f, \operatorname{Im} f (f \text{ 为群同态}) : I, \S 4, \text{ no. } 5.\)

\(\prod_{i \in I} G_i (G_i \text{为群}) : I, \S 4, \text{no. } 8.\)

\(G_{1} \times_{H} G_{2}: I, \S 4, no. 8.\)

\(\prod_{i \in I} G_i (G_i \text{为群}) : I, \S 4, \text{no. } 9.\)

\(x \equiv y \pmod{a}, x \equiv y (a) (a, x, y \text{ 为有理整数}) : I, \S 4, \text{ no. } 10.\)

\(v_{p}(a)\) ( \(p\) 一个素数， \(a\) 一个有理整数）：I，§ 4，第 10 号。

\(\operatorname{Aut}(\mathbf{G}), \operatorname{Int}(\mathbf{G}), \operatorname{Int}(x) (\mathbf{G}\) 一个群， \(x \in \mathbf{G}) : \mathbf{I}, \S 5\) ，第 3 号。

\(N_{G}(A), N(A)\) （G 是一个群，A ⊂ G）：I，§ 5，第 3 号。

\(C_{G}(A), C(A)\) （G 是一个群，A ⊂ G）：I，§ 5，第 3 号。

E/G， \(G\backslash E\) （G 是作用在 E 上的群）：I，§ 5，第 4 号。

\(G\backslash E/H\) （G，H 是通过交换作用作用在 E 上的群）：I，§ 5，第 4 号。 \(S_{n}\) ：I，§ 5，第 7 号。

\(\tau_{x,y}\) （支撑为 \(\{x,y\}\) 的对换）：I，§ 5，第 7 号。

\(\varepsilon (\sigma),\varepsilon_{\sigma}\) （σ 是一个置换）：I，§ 5，第 7 号。

\(\mathfrak{A}_{\mathbf{E}}, \mathfrak{A}_n: \mathbf{I}, \S 5, \text{no. } 7.\)

\(\mathbf{F} \xrightarrow{i} \mathbf{E} \xrightarrow{p} \mathbf{G}\) （E，F，G 是群）：I，§ 6，第 1 号。

\(\mathbf{F} \times_{\tau} \mathbf{G}, \mathcal{E}_{\tau}\) ( \(\tau\) 为 \(\mathbf{G}\) 到 \(\operatorname{Aut}(\mathbf{F})\) 的同态 ): \(\mathbf{I}, \S 6\) ，第 1 号。

\(^{g}f(f\in\mathbf{F},g\in\mathbf{G}):I,\S6,no.1.\)

\((f,g)\cdot_{\tau}(f',g')\) ( \(f,f'\) 在 F 中，g， \(g'\) 在 G 中）：I，§ 6，第 1 号。

\((x,y)\) ，（A，B） \((x,y\) 元素，A，B 为群 G 的子集）：I，§ 6，第 2 号。

D(G)：I，§ 6，第 2 号。

\(C^{n}(G)\) ：I，§ 6，第 3 号。

\(D^{n}(G)\) ：I，§ 6，第 4 号。

\(E^{G}\) （G为作用于E的群）：I，§ 6，第 5 号。

\(M_{n}(X), M(X)\) （X为集合）：I，§ 7，第 1 号。

\(l(w)\) （w 为 \(\mathbf{M}(\mathbf{X})\) 的元素）：I，§7，第1号。

\(ww', w.w' (w, w'\text{为 M(X) 的元素}): I, \S 7, no. 1.\)

\(l(w)\) （w 为 X 上的一个词）：I，§ 7，第 2 号。

\(ww^{\prime}, w.w^{\prime}\) ( \(w, w^{\prime}\) 为 X 上的字)：I，§ 7，第 2 号。

Mo(X)（X 为集合）：I，§ 7，第 2 号。

\(l(\sigma)\) ( \(\sigma\) 为一个分解）：I，§7，第3节。

\(*\mathbf{G}_i,\mathbf{G}_1*\mathbf{G}_2\) ( \(\mathbf{G}_1,\mathbf{G}_2,\mathbf{G}_i\) 群）：I，§ 7，第 3 号。

\(\langle \tau_{1},\ldots ,\tau_{n};r_{1},\ldots ,r_{m}\rangle\) ( \(\tau_{j}\) 生成元， \(r_i\) 关系子）：I，§ 7，第 6 号。

\(\langle \tau_{1},\ldots ,\tau_{n};u_{1} = v_{1},\ldots ,u_{m} = v_{m}\rangle\) ( \(\tau_{j}\) 生成元， \(u_{i},v_{i}\) 自由群的元素）：I，§ 7，第 6 节。

\(\mathbf{Z}^{(\mathbf{x})}, \mathbf{N}^{(\mathbf{x})}\) （X 为集合）：I，§ 7，第 7 节。

0，1（环的元素）：I，§ 8，第 1 节。

\(A^{0}\) （A 为环）：I，§ 8，第 3 节。

\begin{enumerate}
\def\labelenumi{(\alph{enumi})}
\tightlist
\item
  （a 为 A 的元素）：I，§ 8，第 6 节。
\end{enumerate}

\(\sum_{\lambda} a_{\lambda}\) ( \(a_{\lambda}\) 理想）：I，§ 8，第 6 节。

\(x \equiv y (\text{mod } a), x \equiv y (a) (a \text{ 为理想}): I, § 8, no. 7.\)

A/a（a 为双边理想）：I，§ 8，第 7 节。

ab（a，b 为双边理想）：I，§ 8，第 9 节。

\(\mathbf{A}[\hat{\mathbf{S}}^{-1}]\) （S 为环的子集 \(\mathbf{A}\) ）：I，§ 8，第 12 节。

\(\mathbf{F}_p\) ( \(p\) 素数）：I，§ 9，第 1 节。

Q：I，§ 9，第 4 节。

\(Q_{+}\) ：I，§ 9，第 4 节。

\(|x|\) ，sgn x（x 为有理数）：I，§ 9，第 4 节。

in(G)：I，§ 4，习题 13。

\(D_{n}\) ：I，§ 6，习题 4。

Q：I，§ 6，习题 4。

A ≈ B：I，§ 6，习题 39。

\(e(\mathbf{G})\) ：I，§ 7，习题 39。

\(d_{n}(\mathbf{X}) : \mathbf{I}, \S 7, \text{Exercise 39.}\)

\(A_{s}, A_{d}\) （A 为环）：II，§1，第 1 节。

\(\sum_{i \in I} x_i((x_i)_{i \in I}\) 有限支撑模的元素族：II，§ 1，第 1 节。

\(\operatorname{Hom}_{\mathbf{A}}(\mathbf{E}, \mathbf{F})\) , \(\operatorname{Hom}(\mathbf{E}, \mathbf{F})\) , \((\mathbf{E}, \mathbf{F}, \mathbf{A}\) -模）：II，§ 1，第 2 节。

\(\operatorname{End}_{\mathbf{A}}(\mathbf{E}), \operatorname{End}(\mathbf{E}), \operatorname{Aut}(\mathbf{E}), \mathbf{GL}(\mathbf{E})\) （E 为 A-模）：II，§ 1，第 2 节。

\(\operatorname{Hom}_{\mathbf{A}}(u, v)\) , \(\operatorname{Hom}(u, v)\) ( \(u, v\) 线性映射）：II，§ 1，第 2 节。

\(l_{E}\) （E 为模）：II，§1，第 2 节。

0（零模）：II，§ 1，第 3 节。

\(\operatorname{Ker} u, \operatorname{Im} u, \operatorname{Coim} u, \operatorname{Coker} u\) ( \(u\) 线性映射）：II，§ 1，第 3 节。

\(\prod_{i}f_{i}(f_{i}:E_{i}\to F_{i}\text{ 为线性映射}):II,\S 1,\text{ no. }5.\)

\(\bigoplus_{t \in I} E_t, E_p \oplus E_{p+1} \oplus \cdots \oplus E_q ((E_t)_{t \in I}\) A-模族）：II，§ 1，第 6 节。

\(\sum_{i \in I} f_i (f_i: E_i \to F_i \text{ 为线性映射}): II, \S 1, \text{ no. } 6.\)

\(\bigoplus_{i\in I}f_{i},f_{p}\oplus f_{p+1}\oplus\cdots\oplus f_{q}(f_{i}:E_{i}\to F_{i}\text{ 为线性映射}):II,\S 1,no.6.\)

\(\mathbf{E}^{(\mathrm{I})}\) （E 为模）：II，§ 1，第 6 节。

\(\bigoplus_{i \in I} \mathbf{M}_i ((\mathbf{M}_i)_{i \in I}\) 子模族）：II，§ 1，第 7 节。

长 \(_{\mathbf{A}}(\mathbf{M})\) ，长( \(\mathbf{M}\) )（M 为有限长度的 A-模）：II，§ 1，第 10 节。

\(\delta_{st}\) （克罗内克符号）：II，§ 1，第 11 节。

\(\sum_{t \in \mathbf{T}} \xi_t. t\) （T 为集合， \(\xi_t\) 环的元素）：II，§ 1，第 11 节。

Ann(S)，Ann(x)（S 为模的子集，x 为模的元素）：II，§ 1，第 12 节。

\(\rho_{*}(\mathbf{E}), \mathbf{E}_{[\mathbf{B}]} (\mathbf{E}\) A-模， \(\rho: \mathbf{B} \to \mathbf{A}\) 环同态）：II，§ 1，第 13 节。

\(\rho_{*}(u)\) ( \(\rho: \mathbf{B} \to \mathbf{A}\) 环同态， \(u\) A-线性映射）：II，§1，第 13 节。

E*（E 为模）：II，§ 2，第 3 节。

\(\langle x, x^{*} \rangle\) ( \(x\) 左模 E 的元素， \(x^{*}\) 其对偶 E 的元素 \(^{*}\) ）：II，§2，第 3 节。

\(\langle x^{*}, x \rangle\) ( \(x\) 右模 E 的元素， \(x^{*}\) 其对偶 E 的元素 \(^{*}\) ）：II，§ 2，第 3 节。

\(^{t}u\) （u 为线性或半线性映射）：II，§2，第 5 节。

\(\check{u}\) （u 为同构）：II，§ 2，第 5 小节。

\(\mathbf{E} \otimes_{\mathbf{A}} \mathbf{F}, \mathbf{E} \otimes_{\mathbf{A}} \tilde{\mathbf{F}}\) （E 为右 A-模，F 为左 A-模）：II，§ 3，第 1 小节。

\(x \otimes y\) ( \(x \in \mathbf{E}\) （右模）， \(y \in \mathbf{F}\) （左模））：II，§ 3，第 1 小节。

\(u \otimes v (u, v \text{ 为线性映射}) : \mathbf{II}, \S 3, \text{ no. } 2.\)

\(u \otimes v\) ( \(u, v\) 半线性映射）：II，§ 3，第 3 小节。

\(_{s}\) A \(_{d}\) （A 为环）：II，§ 3，第 4 小节。

\(\mathcal{L}_2(\mathrm{E},\mathrm{F};\mathrm{G})\) （E、F、G 为交换环上的模）：II，§ 3，第 5 小节。

\(\bigotimes_{\lambda \in \mathbf{L}} \mathbf{G}_{\lambda}, \bigotimes_{\lambda \in \mathbf{L}} x_{\lambda} ((\mathbf{G}_{\lambda})\) 一族 \(\mathbf{Z}\) -模， \(x_{\lambda} \in \mathbf{G}_{\lambda}\) 为所有 \(\lambda)\) ：II，§ 3，第 9 小节。

\(\bigotimes_{\lambda \in L} v_{\lambda}\) ( \(v_{\lambda}: G_{\lambda} \to G_{\lambda}'\) Z-线性映射）：II，§ 3，第 9 小节。

\(\bigotimes_{(c,p,q)} \mathbf{G}_{\lambda}, \bigotimes_{(c,p,q)} x_{\lambda}, \bigotimes_{(c)} x_{\lambda}\) ：II，§ 3，第 9 小节。

\(\bigotimes_{(c)}v_{\lambda}\) \((v_{\lambda}\) Z-线性映射）：II，§ 3，第 9 小节。

\(E_{1} \otimes_{A_{1}} E_{2} \otimes_{A_{2}} E_{3} \otimes \cdots \otimes_{A_{n-2}} E_{n-1} \otimes_{A_{n-1}} E_{n}: II,\) § 3，第 9 小节。

\(x_{1} \otimes x_{2} \otimes \cdots \otimes x_{n}\) ：II，§ 3，第 9 小节。

\(u_{1} \otimes u_{2} \otimes \cdots \otimes u_{n}\) ( \(u_{i}\) 线性映射）：II，§ 3，第 9 小节。

\(\mathcal{L}_n(\mathrm{E}_1, \ldots, \mathrm{E}_n; \mathrm{G})\) ( \(\mathrm{E}_1, \ldots, \mathrm{E}_n\) , \(\tilde{\mathrm{G}}\) 交换环上的模）：II，§ 3，第 9 小节。

\(\operatorname{Tr}(u)\) ( \(u\) 交换环上的模的自同态）：II，§ 4，第 3 小节。

\(\rho^{*}(\mathbf{E}), \mathbf{E}_{(\mathbf{B})}\) （E 为 A-模， \(\rho: \mathbf{A} \to \mathbf{B}\) 环同态）：II，§ 5，第 1 小节。

\(\rho^{*}(u), u_{(B)} (\rho: A \to B\) 环同态， \(u\) A-模同态）：II，§ 5，第 1 小节。

\(\dim_{\mathbf{K}} \tilde{\mathbf{E}}, \dim \mathbf{E}, [\mathbf{E} : \mathbf{K}] (\mathbf{E}\) K-向量空间）：II，§ 7，第 2 小节。

\(\dim_{A} E, \dim E (E \text{ 为 A-模，且其任意两个基等势): II, \S 7, \text{no. 2.}}\)

余维数 \(_{E}\) F，余维数 F（F 为向量空间 E 的向量子空间）：II，§ 7，第 3 小节。

\(\operatorname{rg}(u)\) ( \(u\) 向量空间的线性映射）：II，§ 7，第 4 小节。

\(\operatorname{rg}(u)\) ( \(u\) 向量空间的张量积的元素）：II，§ 7，第 8 小节。

\(\dim_{K}E\) , \(\dim E\) （E 为域 K 上的仿射空间）：II，§ 9，第 1 小节。

\(a + \mathbf{t},\mathbf{t} + a\) ( \(a\) 点， \(\mathbf{t}\) 仿射空间的平移）：II，§ 9，第 1 小节。

\(b - a\) （a、b 为仿射空间的点）：II，§ 9，第 1 小节。

\(\sum_{c \in I} \lambda_t x_t ((x_t)_{t \in I}\) 仿射空间的一族点， \((\lambda_t)_{t \in I}\) 一族标量，

有限支撑，使得 \(\sum_{i} \lambda_{i} = 1\) 或 \(\sum_{i} \lambda_{i} = 0\) ：II，§ 9，第 1 小节。

\(\mathbf{P}(\mathbf{V}), \Delta(\mathbf{V})\) （V 为向量空间）：II，§ 9，第 5 小节。

\(\mathbf{P}_n(\mathbf{K}), \Delta_n(\mathbf{K})\) （K 为域）：II，§ 9，第 5 小节。

\(\dim_{K} \mathbf{P}(V)\) , \(\dim \mathbf{P}(V)\) （V 为 K-向量空间）：II，§ 9，第 5 小节。

\(\tilde{\mathbf{K}}, \infty\) （K 为域）：II，§ 9，第 9 小节。

\(\mathbf{PGL}(\mathrm{V}), \mathbf{PGL}_n(\mathrm{K}), \mathbf{PGL}(n, \mathrm{K})\) （K 为域，V 为向量空间）：II，§ 9，第 10 小节。

\(^{t}M\) （M 为矩阵）：II，§ 10，第 1 小节。

\(M^{\prime} + M^{\prime \prime}(M^{\prime},M^{\prime \prime}\) 交换群上的矩阵）：II，§ 10，第 2 小节。

\(f(M', M''), M'M'' (M', M'' matrices): \mathbf{II}, \S 10, \text{no. } 2.\)

\(E_{ij}\) （矩阵单位）：II，§ 10，第 3 小节。

\(\sigma(M), M^{\sigma}\) （M 为矩阵， \(\sigma\) 环同态）：II，§ 10，第 3 小节。

\(M(x), \mathbf{x}\) ( \(x\) 有限生成自由模的元素）：II，§ 10，第 4 小节。

\(M(u)\) （u 为自由模到自由模的同态）：II，§ 10，第 4 小节。

\(M(x), M(u)\) （关于直和分解的矩阵）：II，§ 10，第 5 小节。

\(M(u)\) （u 为半线性映射）：II，§ 10，第 6 小节。

\(\mathbf{M}_{n}(\mathbf{A}), I_{n} \ 1_{n} (\mathbf{A} \ a \ ring): \mathbf{II}, \S 10, \ no. 7.\)

\(\mathbf{GL}_n(\mathbf{A}), \mathbf{GL}(n, \mathbf{A})\) （A 为环）：II，§ 10，第 7 小节。

\(^{t}X^{-1}\) （X 为可逆方阵）：II，§ 10，第 7 小节。

\(\operatorname{diag}(a_{i})_{i \in I}, \operatorname{diag}(a_{1}, a_{2}, \ldots, a_{n}) : II, \S 10, no. 7.\)

\(X_{1} \otimes X_{2}\) ( \(X_{1}, X_{2}\) 交换环上的矩阵）：II，§ 10，第 10 小节。

\(\operatorname{Tr}(X)\) ( \(X\) 交换环上的方阵）：II，§ 10，第 11 小节。

\(\operatorname{rg}(X)\) ( \(X\) 域上的矩阵）：II，§ 10，第 12 小节。

\(\deg (x)\) ( \(x\) 分次群的元素）：II，§ 11，第 1 小节。

\(\mathbf{M}(\lambda_{0})\) （M 为分次模， \(\lambda_{0}\) 次数幺半群的元素）：II，§ 11，第 2 小节。

\(\operatorname{Homgr}_{\mathbf{A}}(\mathbf{M}, \mathbf{N})\) （M、N 为分次环 A 上的分次模）：II，§ 11，第 6 小节。

\(\operatorname{Engr}_{\mathbf{A}}(\mathbf{M}), \mathbf{M}^{*\mathbf{g}\mathbf{r}}\) （M 为分次模）：II，§ 11，第 6 小节。

M:N（M、N 为模）：II，§ 1，习题 24。

\(\begin{bmatrix} a & b \\ d & c \end{bmatrix}\) ( \(a, b, c, d\) 射影直线上的点）：II，§ 9，习题 11。

SL(E)（E 为向量空间）：II，§ 10，习题 12。

PSL(E)（E 为向量空间）：II，§ 10，习题 14。

x.y、xy（代数中的乘法）：III，§ 1，第 1 小节。

\(E^{0}\) （E 为代数）：III，§ 1，第 1 小节。

\(\operatorname{Hom}_{\mathbf{A}\text{-alg.}}(\mathbf{E}, \mathbf{F}) (\mathbf{E}, \mathbf{F} \text{ 为 A-代数}): \mathbf{III}, \S 1, \text{ no. } 1.\)

E/b（b 为代数 E 的双边理想）：III，§ 1，第 2 小节。

\(\tilde{\mathbf{E}}\) （E 为代数）：III，§ 1，第 2 小节。

\(\eta_{c}, \eta_{\mathbf{E}}, \eta\) （E 为代数， \(c\) 单位元）：III，§ 1，第 3 小节。

\(\mathbf{T}(u), \mathbf{N}(u): \mathbf{III}, \S 2, \text{no. } 3.\)

\(\bar{u}\) ：III，第2节，第4条。

H：III，第2节，第5条。

\(\operatorname{Lib}_{\mathbf{A}}(\mathbf{I}), \operatorname{Libas}_{\mathbf{A}}(\mathbf{I}), \operatorname{Libasc}_{\mathbf{A}}(\mathbf{I}) : \operatorname{III}, \S 2, \text{no. } 7.\)

\(\mathbf{U}((x_i)_{i \in \mathbf{I}}): \mathbf{III}, \S 2, \text{no. } 8.\)

\(\mathrm{A}[(x_i)_{i \in \mathbf{I}}], \mathrm{A}[x_i]_{i \in \mathbf{I}}: \mathrm{III}, \S 2, \text{no. } 9.\)

\(\mathrm{A}[(\mathrm{X}_i)_{i\in \mathbf{I}}]\) ：III，§ 2，第 9 号。

\(\mathrm{A}[\mathrm{X}_1,\dots ,\mathrm{X}_n]\) ：III，§ 2，第 9 号。

A{[}X{]}，A{[}X, Y{]}，A{[}X, Y, Z{]}：III，§ 2，第 9 号。

\(X^{\nu}\) （v 为多重指标）：III，第2节，第9条。

\(\sum_{s} \xi_{s} e_{s}, \sum_{s} \xi_{s}. s\) （s 为幺半群的元素）：III，第2节，第10条。

\(\mathrm{A}[[\mathrm{X}_i]]_{i \in \mathbf{I}}: \mathrm{III}, \S 2, \mathrm{no.} 11.\)

\(\sum_{v} \alpha_v X^v: III, \S 2, no. 11.\)

\(\omega(u), \omega_{\mathbf{K}}(u)\) ( \(u\) 形式幂级数）：III，第2节，第11条。

\(\bigotimes_{i \in I} E_i, E_1 \otimes_A E_2 \otimes \cdots \otimes_A E_n, E_1 \otimes E_2 \otimes \cdots \otimes E_n (E_i A-algebras): III,\) 第4节，第1条。

\(\mathbf{E}^{\otimes n}\) （E 为代数）：III，第4节，第1条。

\(\bigotimes_{i \in I} E_i\) （I 为无限集， \(E_i\) 代数）：III，第4节，第5条。

\(\bigotimes_{i \in I} u_i, \bigotimes_{i \in I} x_i\) ( \(u_i\) 代数同态， \(x_i\) 元素，I 无限）：III，第4节，第5条。

\(^{8}\bigotimes_{i\in I}E_{i},^{8}\bigotimes_{i\in I}f_{i},^{8}G^{\otimes n}(E_{i},G\text{ 为分次代数， }f_{i}\text{ 为分次代数同态， }\varepsilon\text{ 是一组交换因子): III, § 4, no. 7.}\)

\(^{g}\bigotimes_{i\in I}E_{i},E^{g}\otimes_{A}F,^{g}G^{\otimes n},^{g}\bigotimes_{i\in I}f_{i},f_{1}^{g}\otimes f_{2},^{g}f^{\otimes n}:III,\S4,no.7.\)

\(\otimes^{n}\mathbf{M},\mathsf{T}^{n}(\mathbf{M}),\operatorname {Tens}^{n}(\mathbf{M}),\mathsf{T}_{\mathbf{A}}^{n}(\mathbf{M}),\mathsf{T}(\mathbf{M}),\operatorname {Tens}(\mathbf{M}),\mathsf{T}_{\mathbf{A}}(\mathbf{M})\) （M 为 A-模）：III，第5节，第1条。

\(\mathsf{T}(u), \mathsf{T}^n(u)\) ( \(u\) 线性映射）：III，第5节，第2条。

\(\mathbf{T}_{\mathbf{J}}^{\mathbf{I}}(\mathbf{M}), \mathbf{T}_{q}^{p}(\mathbf{M})\) （M 为模）：III，第5节，第6条。

\(\alpha_{g}^{f}(z), e_{f}^{g}\) ：III，第5节，第6条。

\(\alpha_{\cdot \nu \rho}^{\lambda \dots \mu}\) III，第5节，第6条。

\(c_{j}^{i}\) ：III，第5节，第6条。

\(S(M), S_{A}(M), \text{Sym}(M): III, \S 6, no 1.\)

\(\mathfrak{J}^{\prime}, \mathfrak{J}_{\mathbf{M}}^{\prime}, \mathfrak{J}_{n}^{\prime}: \mathrm{III}, \S 6, \mathrm{no.~1}\) .

\(S^{n}(\mathbf{M}), S^{n}(u), S(u)\) （M 为模，u 为线性映射）：III，第6节，第1和第2条。

\begin{enumerate}
\def\labelenumi{\alph{enumi}.}
\setcounter{enumi}{18}
\tightlist
\item
  z：III，第6节，第3条。
\end{enumerate}

\(e^{\alpha}\) ( \(\alpha\) 多重指标）：III，第6节，第6条。

\(\Lambda (\mathbf{M}),\Lambda_{\mathbf{A}}(\mathbf{M}),\mathrm{Alt}(\mathbf{M}):\mathrm{III},\S 7,\mathrm{no.~1.}\)

\(\mathfrak{J}^{\prime \prime},\mathfrak{J}_{\mathbf{M}}^{\prime \prime},\mathfrak{J}_{n}^{\prime \prime}\) ：III，第7节，第1条。

\(\bigwedge^{n}(\mathbf{M})\) ：III，第7节，第1条。 \(u \wedge v, x_{1} \wedge x_{2} \wedge \cdots \wedge x_{n}\) ：III，第7节，第1条。 \(\bigwedge(u), \bigwedge^{n}(u)\) ( \(u\) 线性映射）：III，第7节，第2条。 \(x_{\mathrm{H}}\) (H 是 \([1, n]\) 的子集)：III，第 7 节，第 3 号。 \(u\) ( \(x_{1}, \ldots, \hat{x}_{j}, \ldots, x_{n}\) ）：III，第7节，第4号。 \(a.z\) ：III，§7，第4号。 \(A_{n}'(\mathbf{M}), A_{n}''(\mathbf{M})\) ：III，§7，第4号。 \(e_{J}\) ：III，§7，第8号。 \(\rho_{J,K}\) ：III，§7，第8号。 \(\det(u)\) ( \(u\) 一个自同态）：III，§ 8，第 1 号。 \(\det(x_{1}, x_{2}, \ldots, x_{n})\) ( \(x_{j}\) 一个 \(n\) -维自由A-模中的向量）：III，§8，第1节。 \(\det(X)\) ( \(X\) 一个矩阵）：III，§8，第3号。 \(\det(\xi_{ij})_{1 \leq i \leq n, 1 \leq j \leq n}, \det(\xi_{ij})\) ：III，§8，第3号。 \(\left| \begin{array}{cccc}\xi_{11} & \cdots & \xi_{1n} \\ . & . & . & . \\ . & . & . & . \\ \xi_{n1} & \cdots & \xi_{nn} \end{array} \right|: \text{III}, \text{§ 8, no. 3.}\) \(X_{\mathrm{H,K}}, X^{jt}\) ( \(X\) 一个矩阵）：III，§8，第5和第6号。 \(\mathbf{SL}_{n}(\mathbf{A}), \mathbf{SL}(n, \mathbf{A})\) ：III，§ 8，第 9 号。 \(p.x\) ( \(p \in A[\mathrm{X}]\) , \(x\) 一个 A-模的元素）：III，§ 8，第 10 号。 \(M_{u}, M[X]\) （M 是一个 A-模， \(u\) 一个自同态）：III，§ 8，第 10 号。 \(\chi_{u}(\mathrm{X})\) ( \(u\) 一个 A-模自同态）：III，§ 8，第 11 号。 \(Tr_{M/K}(a), N_{M/K}(a), Pc_{M/K}(a; X)\) （A 是一个 K-代数，M 是一个 A-模， \(a \in A\) ）：III，§ 9，第 1 号。 \(Tr_{A/K}(a), N_{A/K}(a), Pc_{A/K}(a; X)\) （A 是一个 K-代数， \(a \in A\) ）：III，§ 9，第 3 号。 \(D_{A/K}(x_{1}, \ldots, x_{n})\) ( \(x_{j}\) 一个 K-代数 A 的元素）：III，§ 9，第 5 号。 \([u,v]_{e}: [u,v] (u,v elements of a graded algebra) : III, § 10, no. 4.\) \(P(D), P(d_{1}, \ldots, d_{n}) (P a polynomial, d_{j} derivations) : III, § 10, no. 4.\) \(ad_{e}(a), ad(a) (a a homogeneous element of a graded algebra) : III, § 10, no. 6.\) \(D_{K}(B,F) (B a K-algebra, F a (B,B)-bimodule) : III, § 10, no. 7.\) \(D_{A,p}(B,F), D_{A}(B,F) : III, § 10, no. 7.\) \(D_{s} (s an endomorphism) : III, § 10, no. 9.\) \(\Omega_{K}(A), d_{A/K}(x), dx (x an element of a K-algebra A) : III, § 10, no. 11.\) \(D_{i}P, ∂P/∂X_{i} (P a polynomial) : III, § 10, no. 11.\) \(Ω(u), Ω_{0}(u) (u a ring homomorphism) : III, § 10, no. 12.\) \(Ω_{u} (u a K-algebra homomorphism) : III, § 10, no. 12.\) \(M^{*gr} (M a graded module) : III, § 11.\) \(u.v, u.mv (u,v symmetric multilinear mappings) : III, § 11, no. 2.\) \(u \wedge v (u,v alternating multilinear mappings) : III, § 11, no. 2.\) \(θ_{T}, θ_{S}, θ_{A} : III, § 11, no. 5.\)

\(u \perp x, i(x)\) ：III，§ 11，第 6 号。

\(x \perp u, i'(x): \text{III}, \S 11, \text{no. } 6.\)

\(x \perp u, i(u): \text{III}, \S 11, \text{no. } 7.\)

\(u \perp x, i'(u): \text{III}, \S 11, \text{no. } 7.\)

\(\mathbf{G}_p(\mathbf{E}), \mathbf{G}_{n,p}(\mathbf{K}) : \mathbf{III}, \S 11, \text{no. } 13.\)

\(a(x,y,z)\) ：III，附录，第1号。

ME：III，§2，练习13。

E * F：III，§5，练习6。

R{[}a{]}：III，§6，练习4。

\(\tilde{X}\) ：III，§11，练习9。

K{[}X；σ，d{]}第三卷，第10节，习题3。

